\documentclass[11pt]{article}

\usepackage[margin=1in]{geometry}
\usepackage{amsmath,amssymb,amsthm}
\usepackage{bm}
\usepackage{enumitem}
\usepackage{booktabs}
\usepackage{hyperref}

\newcommand{\E}{\mathbb{E}}
\newcommand{\Prob}{\mathbb{P}}
\newcommand{\given}{\,\vert\,}
\newcommand{\ind}[1]{\mathbf{1}\!\left\{#1\right\}}
\newcommand{\dd}{\,\mathrm{d}}

\title{The Event Projection Problem in Event-Driven Clinical Trials\\
\large A Bayesian Event-and-Time Landmark Formulation (BEATLES)}
\author{}
\date{}

\begin{document}
\maketitle

\section{Problem statement}

Many confirmatory trials---most oncology trials in particular---are
\emph{event-driven}: a milestone such as an interim analysis or the final
database lock (DBL) is triggered not by calendar time but by the accumulation
of a target number of events of interest (e.g., progression or death). At any
interim transfer date the project statistician must answer two dual prediction
questions, conditioning on the data observed so far:

\begin{enumerate}[label=(\textbf{Q\arabic*}),leftmargin=*]
  \item \textbf{Event prediction.} For a fixed future horizon
  $t^{1}$, how many additional events will have accrued by $t^{1}$?
  \item \textbf{Time (landmark) prediction.} For a fixed target number of
  additional events $n_{e}$, how long must the trial run before that target
  is reached, i.e.\ the waiting time $W(n_{e})$?
\end{enumerate}

Both quantities are random because the future event process is unobserved at
the interim. The statistical task is therefore to produce their full
\emph{predictive distributions} (point estimate plus a calibrated prediction
interval), not merely a point forecast, so that downstream operational
timelines can be planned with their uncertainty made explicit.

\section{Data and observation scheme}

Let $i = 1,\dots,n$ index enrolled patients. For each patient there is a latent
event time $T_i$ and a latent dropout (loss-to-follow-up / permanent censoring)
time $C_i$. We assume $T_i \perp C_i$. In addition each patient has a fixed
\emph{administrative censoring} time $a_i$, the time already accrued in the
trial at the interim transfer date. The observed time is
\[
X_i \;=\; \min\{T_i,\,C_i,\,a_i\}.
\]
Define the cause indicators
\[
\Delta^{e}_i = \ind{T_i < \min\{C_i,a_i\}} \;(\equiv c_{i1}),
\qquad
\Delta^{c}_i = \ind{C_i < \min\{T_i,a_i\}} \;(\equiv c_{i2}).
\]
This partitions patients into three observation types at the interim:

\begin{center}
\begin{tabular}{lccl}
\toprule
Type & $\Delta^e_i$ & $\Delta^c_i$ & Meaning \\
\midrule
Event observed        & $1$ & $0$ & $X_i = T_i$ (failure recorded) \\
Permanently censored  & $0$ & $1$ & $X_i = C_i$ (dropout / LTFU) \\
Administratively censored & $0$ & $0$ & $X_i = a_i$ (still at risk) \\
\bottomrule
\end{tabular}
\end{center}

The set of administratively censored patients,
$A = \{\, i : \Delta^e_i = \Delta^c_i = 0 \,\}$, is exactly the
\emph{at-risk} cohort whose future outcomes drive the projection; let $n_o$ be
the number of events \emph{already} observed at the interim. Data on patients
who have had an event or dropped out are fixed and will not change once the
trial resumes.

\section{Cause-specific hazards and the likelihood}

Under the independence of $T_i$ and $C_i$, the cause-specific hazard for the
event coincides with its marginal hazard,
\[
\lambda^{e}(t)
= \lim_{\Delta t \to 0}
\frac{\Prob\!\big(T_i \in [t,t+\Delta t),\,\Delta^e_i = 1 \given \min\{T_i,C_i,a_i\}\ge t\big)}{\Delta t}
= \lim_{\Delta t \to 0}
\frac{\Prob\!\big(T_i \in [t,t+\Delta t) \given T_i \ge t\big)}{\Delta t},
\]
and symmetrically for the dropout hazard $\lambda^{c}(t)$. With the standard
relations $S(t) = \exp\{-\int_0^t \lambda(x)\dd x\}$ and $f(t)=\lambda(t)S(t)$,
the contribution of each observation type yields the competing-risks
likelihood (the slides write $f,g$ for the event/dropout densities,
$S_T,S_C$ for the survival functions, and $\bm\lambda_T,\bm\lambda_C$ for the
parameter vectors):
\begin{align}
L(\bm\lambda^{e},\bm\lambda^{c} \given X)
&= \prod_{i=1}^{n}
\Big\{ f^{e}(x_i)\,S^{c}(x_i) \Big\}^{\Delta^e_i}
\Big\{ f^{c}(x_i)\,S^{e}(x_i) \Big\}^{\Delta^c_i}
\Big\{ S^{e}(x_i)\,S^{c}(x_i) \Big\}^{(1-\Delta^e_i)(1-\Delta^c_i)}
\notag\\[4pt]
&= \prod_{i=1}^{n}
\underbrace{\lambda^{e}(x_i)^{\Delta^e_i} S^{e}(x_i)}_{\text{event process}}
\;\times\;
\underbrace{\lambda^{c}(x_i)^{\Delta^c_i} S^{c}(x_i)}_{\text{dropout process}} .
\label{eq:lik}
\end{align}
Because the event and dropout factors separate and share no parameters, the
two hazard processes can be fit independently.

\section{The piecewise-exponential model}

Each hazard is modeled as piecewise constant on $N$ intervals delimited by
change points $T_1 < T_2 < \dots < T_{N+1}$:
\[
h(t) \;=\; \sum_{k=1}^{N} \lambda_k \, \ind{T_k \le t < T_{k+1}}.
\]
In BEATLES the number of pieces $N$ is user-specified, and the change points
$T_k$ default to equally spaced empirical quantiles of the observed event times
(e.g.\ $N=5$ places knots at the $20,40,60,80\%$ quantiles). A separate
piecewise-exponential model is used for the event and dropout hazards, giving
parameter vectors $\bm\lambda^{e}=(\lambda^e_1,\dots,\lambda^e_N)$ and
$\bm\lambda^{c}$. (Weibull and log-normal alternatives are also available to
relax the constant-rate-in-final-piece restriction.)

\section{Bayesian estimation via the Poisson equivalence}

Maximizing or sampling \eqref{eq:lik} under the piecewise-exponential hazard is
equivalent to a Poisson regression. Let $j$ index the piece, let
$c_{ij} = c_i \, \ind{T_j \le t_i < T_{j+1}}$ be the event indicator for
patient $i$ falling in piece $j$, and let $t_{ij}$ be the time patient $i$ is
exposed within piece $j$. Then
\[
c_{ij} \;\sim\; \mathrm{Poisson}\!\big(t_{ij}\,\lambda_j\big),
\]
with conjugate gamma priors on the piece rates,
\[
\lambda_k \;\sim\; \mathrm{Gamma}(0.01,\,0.01), \qquad k = 1,\dots,N .
\]
The posterior is
\[
p(\bm\lambda^{e},\bm\lambda^{c} \given X)
\;\propto\;
L(\bm\lambda^{e},\bm\lambda^{c}\given X)\,
\pi(\bm\lambda^{e})\,\pi(\bm\lambda^{c}),
\]
and, by the factorization above, posterior draws of $\bm\lambda^{e}$ and
$\bm\lambda^{c}$ are obtained separately by MCMC (the implementation uses JAGS
/ \texttt{rjags}; a core re-implementation in C\texttt{++} gives a $>\!50\times$
speed-up). The result is a set of posterior samples
$\{\bm\lambda^{(j)}\}_{j=1}^{M}$ that carry the parameter uncertainty into all
predictions.

\section{Prediction: posterior predictive distribution}

Future outcomes are generated only for patients who are unobserved at the
interim: the at-risk set $A$ ($\Delta^e=\Delta^c=0$) and, optionally, a
\emph{future enrollment} cohort specified by an enrollment rate and duration.
For an unobserved time $\tilde{y}$ the predictive law marginalizes over the
posterior,
\[
p(\tilde{y},\bm\lambda \given y) \;=\; p(\tilde{y}\given \bm\lambda, y)\,
p(\bm\lambda \given y),
\]
so each posterior draw $\bm\lambda^{(j)}$ produces one posterior predictive
realization of the future event stream.

\subsection*{Event prediction (Q1)}

For an at-risk patient $i$ who has already survived to $t^{0}_i$, the
probability of experiencing the event within the window $(t^{0}_i, t^{1}_i]$
\emph{before} dropping out is
\begin{align}
p_i\!\left(t^{0}_i, t^{1}_i\right)
&= \Prob\!\left(T_i \le t^1_i,\; T_i < C_i \;\middle|\; T_i > t^0_i,\, C_i > t^0_i\right)
\notag\\[3pt]
&= \frac{\displaystyle\int_{t^{0}_i}^{t^{1}_i}
\lambda^{e}(t)\,S^{e}(t)\,S^{c}(t)\dd t}
{\,S^{e}(t^{0}_i)\,S^{c}(t^{0}_i)\,}.
\label{eq:eventprob}
\end{align}
Summing over the at-risk set and adding the events already banked gives the
expected total and a sampled total at horizon $t^1$:
\[
\E\{N(I)\} = n_o + \sum_{i \in A} p_i\!\left(t^0_i, t^1_i\right),
\qquad
\widetilde{N}(I) = n_o + \sum_{i \in A}\widetilde{N}_i\!\left(t^0_i, t^1_i\right),
\]
where each $\widetilde{N}_i \sim \mathrm{Bernoulli}\big(p_i(t^0_i,t^1_i)\big)$.
Repeating across posterior draws yields the predictive distribution of the
event count. (This per-patient Bernoulli simulation is the computationally
expensive branch: $\approx 10$ min per $1000$ draws.)

\subsection*{Time / landmark prediction (Q2)}

The waiting time to the $n_e$-th additional event is obtained by direct
forward simulation of the future trajectory for each posterior draw:

\begin{enumerate}[leftmargin=*]
  \item Select the $j$-th posterior sample $\big(\bm\lambda^{e,(j)},
  \bm\lambda^{c,(j)}\big)$.
  \item For every at-risk patient $i \in A$, draw a future event time and
  dropout time from the \emph{conditional} (truncated) piecewise-exponential,
  i.e.\ $\widetilde{T}_i \sim f^{e}(\cdot \given \bm\lambda^{e,(j)})$
  restricted to $\widetilde{T}_i > t^0_i$, and likewise
  $\widetilde{C}_i > t^0_i$. Set
  $\widetilde{X}_i = \min\{\widetilde{T}_i,\widetilde{C}_i\}$ and
  $\widetilde{\Delta}^e_i = \ind{\widetilde{T}_i < \widetilde{C}_i}$.
  \item Order the simulated event times
  $\widetilde{X}_{(1)} \le \dots \le \widetilde{X}_{(k)}$ among patients with
  $\widetilde{\Delta}^e_i = 1$.
  \item If $k < n_e$ set $W^{(j)}(n_e) = \infty$; otherwise
  $W^{(j)}(n_e) = \widetilde{X}_{(n_e)}$ (reported as additional time
  $=$ event time $-$ time already in trial).
  \item Repeat over all posterior draws to obtain the predictive distribution
  of $W(n_e)$.
\end{enumerate}
This branch is fast ($\approx 1$ s per $1000$ draws). Predictions are summarized
by the posterior predictive median and a $(1-\alpha)$ credible/prediction
interval (e.g.\ the $\alpha/2,\,25,\,50,\,75,\,(1-\alpha/2)$ percentiles).

\section{Practical considerations}

\paragraph{Choice of pieces.} Simulation evidence indicates that the number and
placement of pieces has only a minor effect on the projections; when in doubt
it is safer to use more pieces than fewer, and to fit several models to gauge
sensitivity rather than to underfit or overfit the hazard.

\paragraph{Overestimation of events.} A recurring operational failure mode is
that the model predicts \emph{more} events at the data-transfer date than were
actually observed (e.g.\ predicted $146$ vs.\ actual $107$), producing overly
aggressive timelines. Documented mitigations include: (i) resetting the last
observed date of at-risk subjects to the transfer date (assumes no events occur
between last contact and transfer, pushing the projection latest); (ii) adding
the overestimated count to the target number of events; and (iii) anchoring
projected event times to the next scheduled scan/visit date. Across
retrospective oncology studies the standard analysis underestimated the DBL cut
date by $7$--$14$ weeks, whereas the overestimation-corrected analyses landed
within about $3$ weeks.

\paragraph{Communication.} Because event rates can accelerate or slow, results
are best reported as a \emph{range} of likely dates with the sources of
variability stated explicitly: the within-projection predictive uncertainty,
and the between-method variability across parametric models and sensitivity
analyses.

\end{document}